\section{Symbolic pressure-coefficient audit}
\label{sec:symbolic-pressure-audit}

The symbolic audit uses the general zeroth-order pressure scale
\[
  E_p^{(0)} = N_p \frac{4\pi}{3}\mathcal L_K .
\]
The target value \(E_p^{(0)}=(16\pi/3)\mathcal L_K\) is obtained iff
\[
  N_p = 4.
\]
Thus the coefficient \(16\pi/3\) is derived only after deriving the
four-sector pressure count \(N_p=4\).

For the shell correction, with
\[
  \eta_K = \frac{1}{2\chi_R\mathcal L_K},
\]
write
\[
  E_p^{\rm shell}
  =
  E_p^{(0)}\left(1-\sigma\eta_K^2\right)
  =
  E_p^{(0)}\left[
  1-\frac{\sigma}{4\chi_R^2\mathcal L_K^2}
  \right].
\]
At \(\chi_R=2\), the target coefficient \(11/48\) requires
\[
  \sigma = \frac{11}{3}.
\]
The proposed second-variation decomposition is
\[
  \sigma = 3 + \frac{2}{3} = \frac{11}{3},
\]
where the \(3\) term is assigned to the spherical volume second variation and
the \(2/3\) term to isotropic transverse averaging.  This is a derivation target:
the decomposition must be obtained from a controlled shell/NLS calculation to
upgrade the coefficient from closure hypothesis to derived result.

The radius closure is represented by the minimal action
\[
  A_\chi=\chi_R+\frac{\lambda_\chi}{\chi_R},
\]
whose stationary point is
\[
  \chi_R=\sqrt{\lambda_\chi}.
\]
Therefore \(\chi_R=2\) requires
\[
  \lambda_\chi = 4.
\]

\begin{table}[h]
\centering
\caption{Symbolic gates for the pressure-cell derived label.}
\begin{tabular}{@{}llll@{}}
\toprule
Gate & Target & Sufficient assumption & Status\\
\midrule
pressure volume prefactor & 16*pi/3 & N_p=4 pressure sectors & passed_if_Np_derived \\
finite-shell correction & 11/48 & chi_R=2 and sigma=3+2/3=11/3 & passed_if_sigma_decomposition_derived \\
cell-radius closure & chi_R=2 & lambda_chi=4 & passed_if_lambda_chi_derived \\
phase stiffness normalization & K_cell=E_eff/(8*pi) & q_phi=1 canonical unit phase normalization & passed_if_phase_hessian_derived \\
\bottomrule
\end{tabular}
\end{table}
